\documentclass[12pt,a4paper]{article}
\usepackage[T1]{fontenc}
\usepackage{lmodern}
\usepackage[utf8]{inputenc}
\usepackage[ngerman,latin]{babel}
\usepackage{amsmath,amssymb}
\usepackage{booktabs}
\usepackage{textcomp}
\usepackage{geometry}
\usepackage{fancyhdr}

\geometry{margin=2.5cm}

\pagestyle{fancy}
\fancyhf{}
\fancyhead[C]{\thepage}
\renewcommand{\headrulewidth}{0pt}

\title{Carl Friedrich Gau\ss{} Werke --- Band 6}
\author{Carl Friedrich Gau\ss{}}
\date{}

\begin{document}

\maketitle
\thispagestyle{empty}

\subsection*{8.}

Ad maiorem illustrationem praeceptorum praecedentium, calculum integrum,
per quem systema elementorum tertium ex oppositionibus annorum 1805, 1807,
1808, 1809 determinatum est, hic apponam. Supposui longitudinem nodi ascen-
dentis pro initio anni 1803, $172\textdegree\ 28'\ 46\textquotedbl.8 = \beta$, inclinationem orbitae $34\textdegree\ 37'\ 31\textquotedbl.5$;
reduxi illam ad tempora singularum oppositionum, addendo praecessionem
$2'\ 26\textquotedbl.01$, $3'\ 37\textquotedbl.50$, $4'\ 39\textquotedbl.12$, $5'\ 37\textquotedbl.11$. Hinc derivavi ex longitudinibus
latitudines heliocentricas
\begin{align*}
\mathfrak{g} &= -33\textdegree\ 40'\ 50\textquotedbl.63 \\
\mathfrak{g}' &= +28\textdegree\ 14'\ 51\textquotedbl.24 \\
\mathfrak{g}'' &= +27\textdegree\ 20'\ 55\textquotedbl.86 \\
\mathfrak{g}''' &= -4\textdegree\ 52'\ 28\textquotedbl.44
\end{align*}

Porro prodeunt coefficientes
\begin{align*}
c &= +0.1252, \quad f = -0.9870 \\
c' &= -0.3366, \quad f' = +0.8917 \\
c'' &= +0.3609, \quad f'' = +0.8727 \\
c''' &= +0.6803, \quad f''' = -0.1811
\end{align*}

Porro prodeunt argumenta latitudinum
\begin{align*}
&257\textdegree\ 25'\ 6\textquotedbl.73 \\
&56\textdegree\ 24'\ 5\textquotedbl.96 \\
&126\textdegree\ 2'\ 55\textquotedbl.07 \\
&188\textdegree\ 36'\ 2\textquotedbl.69
\end{align*}
e quibus elementa elliptica elicere oportet. Ad hunc finem statuo angulum
$\varphi = 14\textdegree\ 10'\ 4\textquotedbl.08$ (uti in systemate elementorum secundo inventus erat), longi-
tudinemque perihelii, pro initio anni 1803, $= 121\textdegree\ 5'\ 22\textquotedbl.1$. Perihelium (per-
inde ut nodus) respectu stellarum fixarum immobile supponitur, adeoque ipsius
distantia a nodo ascendente constans $= 308\textdegree\ 36'\ 35\textquotedbl.3$. Hinc habemus per for-
mulas notas:

\begin{center}
\begin{tabular}{@{}l@{\hspace{3em}}l@{}}
$v = 308\textdegree\ 47'\ 30\textquotedbl.66$, & $M = 327\textdegree\ 54'\ 28\textquotedbl.87$, \\
$v' = 107\textdegree\ 55'\ 26\textquotedbl.71$, & $M' = 79\textdegree\ 49'\ 5\textquotedbl.20$, \\
$v'' = 177\textdegree\ 15'\ 13\textquotedbl.11$, & $M'' = 175\textdegree\ 54'\ 19\textquotedbl.77$, \\
$v''' = 239\textdegree\ 59'\ 27\textquotedbl.39$, & $M''' = 266\textdegree\ 29'\ 57\textquotedbl.59$
\end{tabular}
\end{center}

Porro
\begin{align*}
m &= -0.68517, \quad n = +1.18569 \\
m' &= -1.06482, \quad n' = -2.01318 \\
m'' &= -1.59701, \quad n'' = -0.12921 \\
m''' &= -1.18352, \quad n''' = +1.93947
\end{align*}

Denotando igitur correctionem longitudinis perihelii per $d\varpi$, atque correctionem
anguli $\varphi$ per $d\varphi$, erit motus medius sidereus

ab oppositione prima usque ad secundam intra dies 521.140706
\[111\textdegree\ 30'\ 41\textquotedbl.97 - 0.37965\,d\varpi - 3.19887\,d\varphi\]

ab oppositione secunda usque ad tertiam intra dies 449.277674
\[96\textdegree\ 8'\ 1\textquotedbl.82 - 0.53219\,d\varpi + 1.88397\,d\varphi\]

ab oppositione tertia usque ad quartam intra dies 422.786667
\[90\textdegree\ 35'\ 28\textquotedbl.72 + 0.41849\,d\varpi + 2.06868\,d\varphi\]

Hinc eruuntur tres expressiones pro valore medio diurno sidereo
\begin{align*}
770\textquotedbl.31398 &- 0.0007285\,d\varpi - 0.0061382\,d\varphi \\
770\textquotedbl.30718 &- 0.0011845\,d\varpi + 0.0041933\,d\varphi \\
771\textquotedbl.37892 &+ 0.0009780\,d\varpi + 0.0048930\,d\varphi
\end{align*}
e quarum comparatione demanant aequationes:
\begin{align*}
0 &= 0\textquotedbl.00680 + 0.0004560\,d\varpi - 0.0103315\,d\varphi \\
0 &= -1\textquotedbl.07174 - 0.0021625\,d\varpi - 0.0006997\,d\varphi
\end{align*}
atque hinc $d\varpi = -488\textquotedbl.82$, $d\varphi = -20\textquotedbl.92$, ac proin

\begin{center}
longitudo perihelii correcta pro initio anni 1803, $120\textdegree\ 57'\ 13\textquotedbl.28$, \\
valorque correctus anguli $\varphi = 14\textdegree\ 9'\ 43\textquotedbl.16$
\end{center}

Porro fit motus medius diurnus sidereus $770\textquotedbl.7985$

anomalia media correcta pro oppositione prima\quad $328\textdegree\ 20'\ 55\textquotedbl.20$

longitudo media pro eadem epocha\quad $89\textdegree\ 20'\ 34\textquotedbl.49$

Denique derivatur ex motu medio sidereo logarithmus semiaxis maioris
$0.4420439$, adeoque logarithmus semiparametri $0.4152361$. Iam quum habea-
mus valores correctos anomaliarum verarum
\begin{align*}
v &= 308\textdegree\ 56'\ 40\textquotedbl.25 \\
v' &= 107\textdegree\ 55'\ 39\textquotedbl.48 \\
v'' &= 177\textdegree\ 34'\ 28\textquotedbl.59 \\
v''' &= 240\textdegree\ 7'\ 36\textquotedbl.21
\end{align*}
computantur logarithmi radiorum vectorum
\begin{align*}
\log r &= 0.3531088 \\
\log r' &= 0.4492406 \\
\log r'' &= 0.5369700 \\
\log r''' &= 0.4716739
\end{align*}

E tabulis solaribus porro habemus
\begin{align*}
\log R &= 9.9937332 \\
\log R' &= 0.0039862 \\
\log R'' &= 0.0065917 \\
\log R''' &= 0.0011160
\end{align*}
unde tandem deducimus
\begin{align*}
\gamma &= -33\textdegree\ 39'\ 48\textquotedbl.15 \\
\gamma' &= +28\textdegree\ 15'\ 0\textquotedbl.73 \\
\gamma'' &= +27\textdegree\ 20'\ 9\textquotedbl.07 \\
\gamma''' &= -4\textdegree\ 52'\ 53\textquotedbl.99
\end{align*}

Ut eruantur coefficientes $a$, $b$ etc., formo hypothesin secundam, retinendo
inclinationem sed augendo longitudinem nodi uno minuto primo, ut sit pro initio
anni 1803, $172\textdegree\ 29'\ 46\textquotedbl.8$. Statuendo dein, ut in hypothesi prima, longitudinem
perihelii pro eadem epocha $= 121\textdegree\ 5'\ 22\textquotedbl.1$, inveniuntur anomaliae verae
\begin{align*}
v &= 308\textdegree\ 48'\ 40\textquotedbl.93 \\
v' &= 107\textdegree\ 47'\ 34\textquotedbl.06 \\
v'' &= 177\textdegree\ 26'\ 22\textquotedbl.25 \\
v''' &= 239\textdegree\ 59'\ 15\textquotedbl.00
\end{align*}
atque hinc (statuendo quoque $\varphi = 14\textdegree\ 10'\ 4\textquotedbl.08$) anomaliae mediae

\begin{center}
\begin{tabular}{@{}l@{\hspace{3em}}l@{}}
$M = 328\textdegree\ 15'\ 51\textquotedbl.59$, \\
$M' = 79\textdegree\ 46'\ 30\textquotedbl.67$, \\
$M'' = 175\textdegree\ 54'\ 32\textquotedbl.83$, \\
$M''' = 266\textdegree\ 29'\ 42\textquotedbl.93$
\end{tabular}
\end{center}

(Ceterum ad hunc computum non opus est, methodum vulgarem adhibere, sed
sufficit, valoribus anomaliarum mediarum in hypothesi prima inventis addere pro-
ducta e coefficientibus $m$, $m'$, $m''$, $m'''$ positive sumtis in differentias inter valores
respectivos anomaliarum verarum in utraque hypothesi, puta
\begin{align*}
+9\textquotedbl.50 \times 0.68517 &= +6\textquotedbl.51 \\
+3\textquotedbl.40 \times 1.06482 &= +3\textquotedbl.62 \\
+2\textquotedbl.48 \times 1.59701 &= +3\textquotedbl.96 \\
-12\textquotedbl.39 \times 1.18352 &= -14\textquotedbl.66)
\end{align*}

Hinc retinendo valores coefficientium $m$, $n$ etc.\ eruitur
\begin{center}
$d\varpi = -468\textquotedbl.21$, longitudo perihelii initio anni 1803, $120\textdegree\ 57'\ 33\textquotedbl.89$ \\
$d\varphi = -20\textquotedbl.62$, $\varphi = 14\textdegree\ 9'\ 43\textquotedbl.46$ \\
motus medius diurnus sidereus $= 770\textquotedbl.7761$ \\
logarithmus semiaxis maioris $= 0.4420523$ \\
longitudo media in oppositione prima $= 89\textdegree\ 20'\ 47\textquotedbl.85$
\end{center}

\begin{align*}
\gamma &= -33\textdegree\ 39'\ 51\textquotedbl.10 \\
\gamma' &= +28\textdegree\ 15'\ 1\textquotedbl.27 \\
\gamma'' &= +27\textdegree\ 20'\ 10\textquotedbl.97 \\
\gamma''' &= -4\textdegree\ 52'\ 54\textquotedbl.66
\end{align*}

Tandem formo hypothesin tertiam statuendo $\beta = 172\textdegree\ 28'\ 46\textquotedbl.6$,
$i = 34\textdegree\ 38'\ 31\textquotedbl.5$, unde perinde ut ante deducitur

\begin{center}
longitudo perihelii pro initio anni 1803, $120\textdegree\ 55'\ 34\textquotedbl.46$ \\
$\varphi = 14\textdegree\ 9'\ 52\textquotedbl.63$ \\
motus medius diurnus sidereus $= 770\textquotedbl.8398$ \\
logarithmus semiaxis maioris $= 0.4420283$ \\
longitudo media in oppositione prima $= 89\textdegree\ 20'\ 20\textquotedbl.65$
\end{center}

\begin{align*}
\gamma &= -33\textdegree\ 39'\ 35\textquotedbl.63 \\
\gamma' &= +28\textdegree\ 15'\ 5\textquotedbl.20 \\
\gamma'' &= +27\textdegree\ 20'\ 9\textquotedbl.32 \\
\gamma''' &= -4\textdegree\ 52'\ 52\textquotedbl.65
\end{align*}

Comparatio valorum latitudinum heliocentricarum $\gamma$, $\gamma'$, $\gamma''$, $\gamma'''$ in tribus
hypothesibus inventarum producit
\begin{align*}
a &= -0.0492, \quad b = +0.2087 \\
a' &= +0.0212, \quad b' = +0.0745 \\
a'' &= +0.0317, \quad b'' = +0.0042 \\
a''' &= -0.0112, \quad b''' = +0.0223
\end{align*}

Habentur itaque quatuor aequationes
\begin{align*}
+62\textquotedbl.48 - 0.1744\,d\beta + 1.1957\,di &= 0 \\
+9\textquotedbl.49 + 0.3578\,d\beta - 0.8172\,di &= 0 \\
-46\textquotedbl.79 - 0.3292\,d\beta - 0.8685\,di &= 0 \\
-25\textquotedbl.55 - 0.6915\,d\beta + 0.2034\,di &= 0
\end{align*}
e quibus per principium quadratorum minimorum deducitur
\begin{align*}
d\beta &= -54\textquotedbl.41 \\
di &= -42\textquotedbl.06
\end{align*}
ita ut habeatur

\begin{center}
longitudo nodi ascendentis pro initio anni 1803, $172\textdegree\ 27'\ 52\textquotedbl.39$ \\
inclinatio orbitae $34\textdegree\ 36'\ 49\textquotedbl.44$
\end{center}

Elementa reliqua vel per comparationem eorum, quae in singulis tribus hy-
pothesibus prodierunt, erui, vel quod accuratius est, per calculum novum longi-
tudinum in orbita atque repetitionem operationum in art.\ 5 explicatarum deter-
inari poterunt. Methodus prima suppeditat

\begin{center}
\begin{tabular}{@{}l@{\hspace{2em}}l@{}}
longitudinem perihelii 1803 & $120\textdegree\ 58'\ 3\textquotedbl.86$ \\
angulum $\varphi$ & $14\textdegree\ 9'\ 36\textquotedbl.25$ \\
longitudinem mediam in oppositione prima & $89\textdegree\ 20'\ 32\textquotedbl.08$ \\
motum medium diurnum sidereum & $770\textquotedbl.7899$ \\
logarithmum semiaxis maioris & $0.4420471$
\end{tabular}
\end{center}

Per methodum alteram invenitur

\begin{center}
\begin{tabular}{@{}l@{\hspace{2em}}l@{}}
longitudo perihelii 1803 & $120\textdegree\ 58'\ 4\textquotedbl.81$ \\
angulus $\varphi$ & $14\textdegree\ 9'\ 36\textquotedbl.63$ \\
longitudo media in oppositione prima & $89\textdegree\ 20'\ 31\textquotedbl.81$ \\
motus medius diurnus sidereus & $770\textquotedbl.7893$ \\
logarithmus semiaxis maioris & $0.4420473$
\end{tabular}
\end{center}

Cum his elementis conveniunt ea, quae supra (art.\ 3) tradidi.

\subsection*{9.}

Quantumvis magnae sint perturbationes, quas Pallas a reliquis planetis pa-
titur, tamen, experientia teste, elementa elliptica quatuor oppositionibus adap-
tata planetae motui intra integrum hoc tempus satis bene satisfaciunt: quin adeo
etiam a motu antecedente ac consequente, nisi intervallum temporis nimium as-
sumatur, parum differunt, ita ut e.\ g.\ elementa secunda in art.\ 3 tradita, in op-
positione anni 1803 tribus, in oppositione anni 1809 quinque minutis primis a
longitudine heliocentrica observata discrepaverint. Quae quum ita sint, ad con-
struendam ephemeridem pro motu planetae futuro semper commodissimum mihi
videtur elementis pure ellipticis uti, quae e quatuor oppositionibus proxime an-
tecedentibus derivata erant, siquidem multitudo aequationum a perturbationibus
oriundarum tanta certo evasura est, ut computus vel unius loci heliocentrici pla-
netae tantum non aeque operosus evadere debeat, ac calculus elementorum el-
lipticorum per praecepta supra tradita. Ita locum planetae geocentricum salvo
errore paucorum minutorum primorum tuto semper praedicere licebit, quae prae-
cisio ad inveniendum planetam sufficit.

\subsection*{10.}

Attamen postulat scientiae dignitas, ut consensui stabiliori prospiciatur,
quem antequam perturbationes in calculum introducantur obtineri non posse ma-
nifestum est. Praematurus fuisset, meo quidem iudicio, calculus tam prolixus
taediique plenus, quamdiu observationum copia tempus nimis exiguum com-
plectebatur, perturbationesque a planetis reliquis oriundae vix se manifestabant.
Nunc vero, ubi motus ellipticus non amplius sufficit, ad omnia loca observata in-
ter se concilianda, tempus adesse videtur, ubi de theoria accuratiori cogitari pot-
est. Quo pacto calculum perturbationum, quas Pallas patitur praesertim a Iove,
commodissime atque exactissime absolvendum censeam, quum methodis pro aliis
planetis adhibitis propter nimiam excentricitatem atque inclinationem vix ac ne
vix quidem uti liceat, mox alio loco fusius explicabo: de elementis autem ellipti-
cis, quae maxime idonea videantur, ut calculus perturbationum ipsis superstrua-
tur, in sequentibus adhuc agam. Eruenda scilicet mihi propono elementa el-
liptica, quae non his vel illis oppositionibus exacte, sed omnibus, quae hactenus
observatae sunt, quam proxime satisfaciant. Methodum quidem, per quam tale
negotium exsequi liceat, iam in Theoria Motus Corporum Coelestium art.\ 187 suc-
cincte descripsi: sed quum non solum ea, quae illic generaliter tractavi, in casu
speciali, ubi loca observata sunt oppositiones, quaedam compendia admittant, sed
etiam quaedam artificia practica, per quae applicationem methodi quadratorum
minimorum faciliorem reddere iamdudum solitus sum, in opere illo desiderentur,
astronomis haud ingratum fore spero, si hosce calculos aliquanto fusius hie tra-
didero. Quum totum negotium versetur in determinatione correctionum elementis
approximatis, quae ab omnibus locis observatis haud multum aberrare supponun-
tur, adiiciendarum, totus labor a duobus momentis pendebit: primo scilicet for-
mari debent aequationes lineares, quas singula loca observata suppeditant, dein
ex his aequationibus valores incognitarum maxime idonei sunt eruendi.

\subsection*{11.}

Sit secundum elementa approximata
\begin{center}
\begin{tabular}{@{}l@{\hspace{2em}}l@{}}
$L$ & longitudo media planetae pro epocha arbitraria \\
$t$ & numerus dierum inde ab epocha usque ad momentum observationis elapsorum \\
$\gamma$ & motus medius diurnus sidereus in minutis secundis \\
$\varpi$ & longitudo perihelii \\
$e = \sin\varphi$ & excentricitas \\
$a$ & semiaxis maior \\
$r$ & radius vector \\
$v$ & anomalia vera \\
$E$ & anomalia excentrica \\
$\beta$ & longitudo nodi ascendentis \\
$i$ & inclinatio orbitae \\
$u$ & argumentum latitudinis \\
$\lambda$ & longitudo heliocentrica \\
$Y$ & latitudo heliocentrica \\
$\delta$ & latitudo geocentrica \\
$R$ & distantia terrae a Sole
\end{tabular}
\end{center}

Ex observatione autem suppono esse
\begin{center}
\begin{tabular}{@{}l@{\hspace{2em}}l@{}}
$\alpha$ & longitudinem heliocentricam \\
$\delta$ & latitudinem geocentricam
\end{tabular}
\end{center}

Denique per $dL$, $d\gamma$, $d\varpi$ etc.\ denoto correctiones quantitatum $L$, $\gamma$, $\varpi$ etc.

Fit itaque $dL + t\,d\gamma$ correctio longitudinis mediae,

$dL + t\,d\gamma - d\varpi$ correctio anomaliae mediae,

adeoque per art.\ 15, 16 Theoriae Motus Corp.\ Coel.
\begin{align*}
dE &= \frac{1}{\cos\varphi}(dL + t\,d\gamma - d\varpi) + \frac{\tan\varphi(2 - e\cos E - e^2\sin^2 E)}{\cos\varphi}\,d\varphi \\
dr &= -da + a\tan\varphi\sin E\,(dL + t\,d\gamma - d\varpi) - a\cos\varphi\cos v\,d\varphi
\end{align*}

Porro fit correctio argumenti latitudinis $du = d\varpi + d\beta - d\beta$,

atque per art.\ 52 Theoriae Motus C.\ C.\ correctio longitudinis heliocentricae:
\[
d\lambda = d\beta - \tan Y\cos(\lambda - \beta)\,di + \frac{r}{\cos Y}\,du
\]

Hinc colligitur
\begin{align*}
d\lambda &= \frac{r}{a\cos Y}\,dL \\
&\quad + \frac{rt\cos\varphi\sin E}{a\cos Y}\,d\gamma \\
&\quad + \Bigl(\frac{r}{a\cos Y} - \frac{\tan Y\cos(\lambda - \beta)}{r\cos Y}\Bigr)\,d\varpi \\
&\quad - \tan Y\cos(\lambda - \beta)\,di \\
&\quad + d\beta + \frac{r}{a\cos Y}\bigl(2 - e\cos E - e^2\bigr)\sin E\,d\varphi
\end{align*}

Porro quum habeatur
\begin{align*}
a^3\gamma^2 &= \text{Const.} \\
r\sin(\delta - \gamma) &= R\sin\delta \\
\tan\gamma &= \tan i\sin(\alpha - \beta)
\end{align*}
fit per differentiationem
\begin{align*}
\frac{da}{a} + \frac{2\,d\gamma}{3\gamma} &= 0 \\
\frac{dR}{R} + \cot(\delta - \gamma)\cdot(d\delta - d\gamma) &= \cot\delta\,d\delta
\end{align*}
sive
\begin{align*}
\frac{dR}{R} &= \frac{\sin\beta\cos(\beta - \gamma)}{\sin\gamma}\,d\delta \\
\frac{d\gamma}{\sin^2\gamma} &= \frac{d\alpha}{\sin^2 i} - \sin 2\gamma\cot(\alpha - \beta)\,d\beta
\end{align*}
unde adhibito valore ipsius $dr$ supra evoluto colligitur
\begin{align*}
d\delta &= -\frac{a\sin\beta\sin(\beta - \gamma)}{r\sin\gamma}\,da \\
&\quad + \frac{a\sin\beta\sin(\beta - \gamma)\tan\varphi\sin E}{r\sin\gamma}\,(dL + t\,d\gamma - d\varpi) \\
&\quad - \frac{a\sin\beta\sin(\beta - \gamma)\cos\varphi\cos v}{r\sin\gamma}\,d\varphi \\
&\quad + \frac{a\sin\beta\sin(\beta - \gamma)}{r\sin\gamma}(2 - e\cos E - e^2)\sin E\,d\varphi \\
&\quad - \frac{\sin\beta\cos(\beta - \gamma)}{\sin\gamma}\,d\delta \\
&\quad - \frac{a\sin\beta\sin(\beta - \gamma)}{r\sin\gamma}\,d\delta
\end{align*}

Hinc valores longitudinis heliocentricae atque latitudinis geocentricae fiunt ex
valoribus correctis elementorum $\lambda + d\lambda$, $\delta + d\delta$, adeoque quaevis oppositio sup-
peditat binas aequationes:
\begin{align*}
\alpha &= \lambda + d\lambda \\
\delta &= \delta + d\delta
\end{align*}

\subsection*{12.}

Applicando haecce praecepta ad sex oppositiones Palladis in art.\ 2 traditas,
si calculum secundo elementorum systemati in art.\ 3 delineato superstruimus, se-
quentes duodecim aequationes obtinemus:

Ex oppositione prima, ubi longitudo computata inventa est $= 277\textdegree\ 36'\ 20\textquotedbl.07$,
latitudo geocentrica $= +4\textdegree\ 26'\ 29\textquotedbl.19$:
\begin{align*}
0 &= -183\textquotedbl.93 + 0.79363\,dL + 143.66\,d\gamma + 0.39493\,d\varpi + 0.95920\,d\varphi \\
&\quad - 0.18856\,d\beta + 0.17387\,di \\
0 &= -6\textquotedbl.81 - 0.02658\,dL + 46.71\,d\gamma + 0.02658\,d\varpi - 0.20858\,d\varphi \\
&\quad + 0.15946\,d\beta + 1.25782\,di
\end{align*}

Ex oppositione secunda, ubi longitudo computata $= 337\textdegree\ 0'\ 36\textquotedbl.04$,
latitudo geocentrica $= +15\textdegree\ 1'\ 46\textquotedbl.71$:
\begin{align*}
0 &= -0\textquotedbl.06 + 0.58880\,dL + 358.12\,d\gamma + 0.26208\,d\varpi - 0.85234\,d\varphi \\
&\quad + 0.14912\,d\beta + 0.17775\,di \\
0 &= -3\textquotedbl.09 + 0.01318\,dL + 28.39\,d\gamma - 0.01318\,d\varpi - 0.07861\,d\varphi \\
&\quad + 0.91704\,d\beta + 0.54365\,di
\end{align*}

Ex oppositione tertia, ubi longitudo computata $= 67\textdegree\ 20'\ 42\textquotedbl.88$,
latitudo geocentrica $= -54\textdegree\ 31'\ 3\textquotedbl.88$:
\begin{align*}
0 &= +0\textquotedbl.02 + 1.73436\,dL + 1846.17\,d\gamma - 0.54603\,d\varpi - 2.05662\,d\varphi \\
&\quad - 0.18833\,d\beta - 0.17445\,di \\
0 &= -8\textquotedbl.98 - 0.12606\,dL - 227.42\,d\gamma + 0.12606\,d\varpi - 0.38939\,d\varphi \\
&\quad + 0.17176\,d\beta - 1.35441\,di
\end{align*}

Ex oppositione quarta, ubi longitudo computata $= 223\textdegree\ 37'\ 25\textquotedbl.39$,
latitudo geocentrica $= +42\textdegree\ 11'\ 28\textquotedbl.07$:
\begin{align*}
0 &= -2\textquotedbl.31 + 0.99584\,dL + 1579.03\,d\gamma + 0.06456\,d\varpi + 1.99545\,d\varphi \\
&\quad - 0.06040\,d\beta - 0.33750\,di \\
0 &= +2\textquotedbl.47 - 0.08089\,dL - 67.22\,d\gamma + 0.08089\,d\varpi - 0.09970\,d\varphi \\
&\quad - 0.46359\,d\beta + 1.22803\,di
\end{align*}

Ex oppositione quinta, ubi longitudo computata $= 304\textdegree\ 2'\ 59\textquotedbl.71$,
latitudo geocentrica $= +37\textdegree\ 44'\ 31\textquotedbl.82$:
\begin{align*}
0 &= +0\textquotedbl.01 + 0.65311\,dL + 1329.09\,d\gamma + 0.38994\,d\varpi - 0.08439\,d\varphi \\
&\quad - 0.04305\,d\beta + 0.34268\,di \\
0 &= +38\textquotedbl.12 - 0.00218\,dL + 38.47\,d\gamma + 0.00218\,d\varpi - 0.18710\,d\varphi \\
&\quad + 0.47301\,d\beta - 1.14371\,di
\end{align*}

Ex oppositione sexta, ubi longitudo computata $= 359\textdegree\ 34'\ 46\textquotedbl.67$,
latitudo geocentrica $= -7\textdegree\ 20'\ 12\textquotedbl.13$:
\begin{align*}
0 &= -317\textquotedbl.73 + 0.69957\,dL + 1719.32\,d\gamma + 0.12913\,d\varpi - 1.38787\,d\varphi \\
&\quad + 0.17130\,d\beta - 0.08360\,di \\
0 &= +117\textquotedbl.97 - 0.01315\,dL - 43.84\,d\gamma + 0.01315\,d\varpi + 0.02929\,d\varphi \\
&\quad + 1.02138\,d\beta - 0.27187\,di
\end{align*}
Sed ex his duodecim aequationibus decimam omnino reiiciemus, quum latitudo
geocentrica observata nimis incerta sit.

\subsection*{13.}

Quum sex incognitas $dL$, $d\gamma$ etc.\ ita determinare non liceat, ut omnibus
undecim aequationibus exacte satisfiat, i.\ e.\ ut singulae incognitarum functiones
quae sunt ad dextram simul fiant $= 0$, valores eos eruemus, per quos functio-
num harum quadrata summam quam minimam efficiant. Facile quidem perspi-
citur, si generaliter functiones lineares incognitarum $p$, $q$, $r$, $s$ etc.\ propositae
sint hae
\begin{align*}
n + ap + bq + cr + ds + \text{etc.} &= w \\
n' + a'p + b'q + c'r + d's + \text{etc.} &= w' \\
n'' + a''p + b''q + c''r + d''s + \text{etc.} &= w'' \\
n''' + a'''p + b'''q + c'''r + d'''s + \text{etc.} &= w'''
\end{align*}
etc., aequationes conditionales, ut $ww + w'w' + w''w'' + w'''w''' + \text{etc.} = \Omega$ fiat
minimum, esse hasce
\begin{align*}
aw + a'w' + a''w'' + a'''w''' + \text{etc.} &= 0 \\
bw + b'w' + b''w'' + b'''w''' + \text{etc.} &= 0 \\
cw + c'w' + c''w'' + c'''w''' + \text{etc.} &= 0 \\
dw + d'w' + d''w'' + d'''w''' + \text{etc.} &= 0
\end{align*}
etc.\ sive, designando brevitatis causa
\begin{align*}
an + a'n' + a''n'' + a'''n''' + \text{etc.} &= [an] \\
aa + a'a' + a''a'' + a'''a''' + \text{etc.} &= [aa] \\
ab + a'b' + a''b'' + a'''b''' + \text{etc.} &= [ab] \quad \text{etc.}
\end{align*}
\begin{align*}
bb + b'b' + b''b'' + b'''b''' + \text{etc.} &= [bb] \\
bc + b'c' + b''c'' + b'''c''' + \text{etc.} &= [bc] \quad \text{etc. etc.}
\end{align*}
per
\begin{align*}
nn + n'n' + n''n'' + n'''n''' + \text{etc.} &= [nn]
\end{align*}
$p$, $q$, $r$, $s$ etc.\ determinari debere per eliminationem ex aequationibus
\begin{align*}
[an] + [aa]p + [ab]q + [ac]r + [ad]s + \text{etc.} &= 0 \\
[bn] + [ab]p + [bb]q + [bc]r + [bd]s + \text{etc.} &= 0 \\
[cn] + [ac]p + [bc]q + [cc]r + [cd]s + \text{etc.} &= 0 \\
[dn] + [ad]p + [bd]q + [cd]r + [dd]s + \text{etc.} &= 0 \quad \text{etc.}
\end{align*}

Attamen quoties multitudo incognitarum $p$, $q$, $r$, $s$ etc.\ paullo maior est, elimi-
natio laborem valde prolixum atque taediosum requirit, quem sequenti modo no-
tabiliter contrahere licet. Praeter coefficientes $[an]$, $[aa]$, $[ab]$ etc.\ (quorum mul-
titudo fit $\frac{1}{4}(n^2 + 3n)$, si multitudo incognitarum $= n$), etiam hunc compu-
tatum suppono $nn + n'n' + n''n'' + n'''n''' + \text{etc.} = [nn]$, perspicieturque facile fieri
\begin{align*}
\Omega &= [nn] + 2[an]p + 2[bn]q + 2[cn]r + 2[dn]s + \text{etc.} \\
&\quad + [aa]pp + 2[ab]pq + 2[ac]pr + 2[ad]ps + \text{etc.} \\
&\quad + [bb]qq + 2[bc]qr + 2[bd]qs + \text{etc.} \\
&\quad + [cc]rr + 2[cd]rs + \text{etc.} \\
&\quad + [dd]ss + \text{etc.} \quad \text{etc.}
\end{align*}

Designando itaque
\[
[an] + [aa]p + [ab]q + [ac]r + [ad]s + \text{etc.} = A
\]
patet, singulas partes ipsius $\frac{\partial\Omega}{\partial p}$, quae factorem $p$ involvunt, in $\Omega$ contineri,
adeoque $\Omega - \frac{A^2}{[aa]}$ esse functionem a $p$ liberam. Quare statuendo
\begin{align*}
[bn] - \frac{[ab]}{[aa]}[an] &= [bn, 1] \\
[cn] - \frac{[ac]}{[aa]}[an] &= [cn, 1] \quad \text{etc.}
\end{align*}
\begin{align*}
[bb] - \frac{[ab]^2}{[aa]} &= [bb, 1] \\
[bc] - \frac{[ab][ac]}{[aa]} &= [bc, 1] \quad \text{etc. etc.}
\end{align*}
erit
\begin{align*}
\Omega - \frac{A^2}{[aa]} &= [nn, 1] + 2[bn, 1]q + 2[cn, 1]r + 2[dn, 1]s + \text{etc.} \\
&\quad + [bb, 1]qq + 2[bc, 1]qr + 2[bd, 1]qs + \text{etc.} \\
&\quad + [cc, 1]rr + 2[cd, 1]rs + \text{etc.} \\
&\quad + [dd, 1]ss + \text{etc.} \quad \text{etc.}
\end{align*}
quam functionem designabimus per $\Omega'$.

\subsection*{14.}

Simili modo ponendo
\[
[bn, 1] + [bb, 1]q + [bc, 1]r + [bd, 1]s + \text{etc.} = B
\]
erit $\Omega' - \frac{B^2}{[bb, 1]}$ functio a $q$ libera, quam statuemus $= \Omega''$. Eodem modo fa-
ciemus
\begin{align*}
[cn, 1] - \frac{[bc, 1]}{[bb, 1]}[bn, 1] &= [cn, 2] \quad \text{etc.}
\end{align*}
\begin{align*}
[cc, 1] - \frac{[bc, 1]^2}{[bb, 1]} &= [cc, 2] \quad \text{etc.}
\end{align*}
etc.\ etc.\ atque
\[
[cn, 2] + [cc, 2]r + [cd, 2]s + \text{etc.} = C
\]
unde $\Omega'' - \frac{C^2}{[cc, 2]}$ functio ab $r$ quoque libera. Eodem modo progrediemur,
usquedum in progressione $\Omega$, $\Omega'$, $\Omega''$ etc.\ ad terminum ab omnibus incognitis li-
berum pervenerimus, qui erit $[nn, n]$, si multitudo incognitarum $p$, $q$, $r$, $s$ etc. denotatur per $n$. Habemus itaque
\[
\Omega = \frac{A^2}{[aa]} + \frac{B^2}{[bb, 1]} + \frac{C^2}{[cc, 2]} + \frac{D^2}{[dd, 3]} + \text{etc.} + [nn, n]
\]

Iam quum $\Omega = ww + w'w' + w''w'' + \text{etc.}$ natura sua valorem negativum obtenere
non possit, facile demonstratur, divisores $[aa]$, $[bb, 1]$, $[cc, 2]$, $[dd, 3]$ etc.\ neces-
sario positivos evadere debere (brevitatis tamen gratia hanc demonstrationem fu-
sius hic non exsequor). Hinc vero sponte sequitur, valorem minimum ipsius $\Omega$
prodire, si fiat $A = 0$, $B = 0$, $C = 0$, $D = 0$ etc.\ Ex his itaque $n$ aequa-
tionibus incognitae $p$, $q$, $r$, $s$ etc.\ determinari debebunt, quod ordine inverso fa-
cillime effici potest, quum manifesto ultima aequatio unicam incognitam impli-
cet, penultima duas et sic porro. Simul haec methodus eo nomine se commen-
dat, quod valor minimus aggregati $\Omega$ sponte inde innotescit, quippe qui mani-
festo est $= [nn, n]$.

\subsection*{15.}

Applicemus iam haecce praecepta ad exemplum nostrum, ubi $p$, $q$, $r$, $s$ etc. sunt $dL$, $d\gamma$, $d\varpi$, $d\varphi$, $d\beta$, $di$. Calculo accurate absoluto hosce valores numericos inveni:

\begin{center}
\begin{tabular}{@{}l@{\hspace{1em}}l@{\hspace{1em}}l@{\hspace{1em}}l@{}}
$[nn] = 148848$ & $[cn, 1] = -119.31$ & $[df, 2] = -0.17265$ \\
& $[ee, 2] = +2.24325$ \\
$[an] = -371.09$ & $[dn, 1] = -125.18$ & \\
$[bn] = -580104$ & $[en, 1] = +72.52$ & $[ef, 2] = -0.37841$ \\
$[cn] = -113.45$ & $[fn, 1] = -43.22$ & $[ff, 2] = +5.61686$ \\
$[dn] = +268.53$ & $[bb, 1] = +2458225$ & \\
$[en] = +94.26$ & $[bc, 1] = +62.13$ \\
$[nn, 2] = +99034$ \\
$[fn] = -31.81$ & $[bd, 1] = -510.58$ & $[dn, 2] = +25.66$ \\
$[aa] = +5.91569$ & $[be, 1] = +213.84$ & $[en, 2] = +74.23$ \\
$[ab] = +7203.91$ & $[bf, 1] = +73.45$ \\
$[ac] = -0.09344$ & $[cc, 1] = +0.71769$ & $[fn, 2] = +2.75$ \\
$[ad] = -2.28516$ & $[cd, 1] = +1.09773$ & $[dd, 2] = +9.29213$ \\
$[ae] = -0.34664$ & $[ce, 1] = -0.05852$ & $[de, 2] = -0.36175$ \\
$[af] = -0.18194$ & $[cf, 1] = +0.26054$ & $[df, 2] = -0.57384$ \\
$[bb] = +10834225$ & $[dd, 1] = +11.12064$ & $[ee, 3] = +2.23754$ \\
$[bc] = -49.06$ & $[de, 1] = -0.50528$ & $[ef, 3] = -0.35532$ \\
$[bd] = -3229.77$ & $[df, 1] = -0.18790$ & $[ff, 3] = +5.52342$ \\
$[be] = -198.64$ & $[ee, 1] = +2.26185$ \\
$[bf] = -143.05$ & $[ef, 1] = -0.37202$ & $[nn, 4] = +98963$ \\
$[cc] = +0.71917$ & $[ff, 1] = +5.61905$ & $[en, 4] = +75.23$ \\
$[cd] = +1.13382$ & & $[fn, 4] = +4.33$ \\
$[ce] = +0.06400$ & $[nn, 2] = +117763$ & $[ee, 4] = +2.22346$ \\
$[cf] = +0.26341$ & $[cn, 2] = -115.81$ & $[ef, 4] = -0.37766$ \\
$[dd] = +12.00340$ & $[dn, 2] = -153.95$ & $[ff, 4] = +5.48798$ \\
$[de] = -0.37137$ & $[en, 2] = +84.57$ \\
$[df] = -0.11762$ & $[fn, 2] = -39.08$ \\
$[ee] = +2.28215$ & $[cc, 2] = +0.71612$ & $[nn, 5] = +96418$ \\
$[ef] = -0.36136$ & $[cd, 2] = +1.11063$ & $[fn, 5] = +17.11$ \\
$[ff] = +5.62456$ & $[ce, 2] = -0.06392$ & $[ff, 5] = +5.42383$ \\
& $[cf, 2] = +0.25868$ \\
& $[dd, 2] = +11.01463$ & $[nn, 6] = +96364$ \\
& $[de, 2] = -0.46088$ & \\
& & $[ff, 6] = +5.41853$
\end{tabular}
\end{center}

Habemus itaque ad determinationem incognitarum sex aequationes sequentes:
\begin{align*}
0 &= +17\textquotedbl.11 + 5.42383\,di \\
0 &= +75\textquotedbl.23 + 2.22346\,d\varphi - 0.37766\,di \\
0 &= +25\textquotedbl.66 + 9.29213\,d\varpi - 0.36175\,d\varphi - 0.57384\,di \\
0 &= -115\textquotedbl.81 + 0.71612\,dn + 1.11063\,d\varphi - 0.06392\,d\beta + 0.25868\,di \\
0 &= -138534 + 2458225\,d\gamma + 62.13\,dn - 510.58\,d\varphi + 213.84\,d\beta \\
&\quad + 73.45\,di \\
0 &= -371\textquotedbl.09 + 5.91569\,dL + 7203.91\,d\gamma - 0.09344\,dn - 2.28516\,d\varpi \\
&\quad - 0.34664\,d\varphi - 0.18194\,di
\end{align*}

unde deducitur
\begin{align*}
di &= -3\textquotedbl.15 \\
d\varphi &= -34\textquotedbl.37 \\
d\varpi &= -4\textquotedbl.29 \\
dn &= +166\textquotedbl.44 \\
d\gamma &= +0\textquotedbl.054335 \\
dL &= -3\textquotedbl.06
\end{align*}

Elementa itaque elliptica correcta, quae omnibus sex oppositionibus quam
proxime satisfaciunt, haec sunt:

\begin{center}
\begin{tabular}{@{}l@{\hspace{2em}}l@{}}
Epocha longitudinis mediae 1803, ad meridianum Gottingensem & $221\textdegree\ 34'\ 53\textquotedbl.64$ \\
Motus medius tropicus diurnus & $770\textquotedbl.5010$ \\
Longitudo perihelii 1803 & $121\textdegree\ 8'\ 8\textquotedbl.54$ \\
Longitudo nodi ascendentis 1803 & $172\textdegree\ 28'\ 12\textquotedbl.43$ \\
Inclinatio orbitae & $34\textdegree\ 37'\ 28\textquotedbl.35$ \\
Excentricitas $= \sin 14\textdegree\ 9'\ 59\textquotedbl.79$ & $0.2447424$ \\
Logarithmus semiaxis maioris & $0.4422071$
\end{tabular}
\end{center}

\subsection*{16.}

Substitutis valoribus correctionum $dL$, $d\gamma$ etc., quos modo invenimus, in
duodecim aequationibus art.\ 12, differentias sequentes inter valores longitudinum
heliocentricarum atque latitudinum geocentricarum observatos atque computatos
obtinemus:

\begin{center}
\begin{tabular}{@{}lcc@{}}
\toprule
& \multicolumn{2}{c}{Differentia} \\
\cmidrule(lr){2-3}
oppositione anni & longitudinis & latitudinis \\
\midrule
1803 & $-68\textquotedbl.6$ & $-371\textquotedbl.00$ \\
1804 & $-19.92$ & $+0.07$ \\
1805 & $+59.18$ & $-19.92$ \\
1807 & $+85.77$ & $+25.01$ \\
1808 & $+216.54$ & --- \\
1809 & $+135.88$ & $-2883.70$ \\
\bottomrule
\end{tabular}
\end{center}

\end{document}
